% !PRINTED 530
\weberpage{530}
%<<CONT
\[
\begin{aligned}
v_1 &= a_1^2 - 4(\alpha+\alpha_1)(\alpha_2+\alpha_3)\\
y_1-y_2 &= (\alpha+\alpha_1)(\alpha_2+\alpha_3) - 2(\alpha\alpha_1+\alpha_2\alpha_3)\\
a_2 &= (\alpha+\alpha_1)(\alpha_2+\alpha_3) + \alpha\alpha_1 + \alpha_2\alpha_3,
\end{aligned}
\]
and thus, if one forms $v_2$, $v_3$ in the same way:
\weq{12}{
\begin{aligned}
3v_1 &= 3a_1^2 - 8a_2 - 4y_1 + 4y_2\\
3v_2 &= 3a_1^2 - 8a_2 - 4y_2 + 4y\\
3v_3 &= 3a_1^2 - 8a_2 - 4y + 4y_1.
\end{aligned}
}
But there is yet another relation, which results from the fact that the product
\[
(\alpha+\alpha_1-\alpha_2-\alpha_3)(\alpha-\alpha_1+\alpha_2-\alpha_3)(\alpha-\alpha_1-\alpha_2+\alpha_3)
\]
is a symmetric function of the $\alpha$, which one easily finds to be equal to
\[
a_1^3 - 4a_1a_2 + 8a_3
\]
We therefore make the signs of the three root quantities depend on one another by the relation
\weq{13}{\sqrt{v_1}\sqrt{v_2}\sqrt{v_3} = a_1^3 - 4a_1a_2 + 8a_3}
and set:
\weq{14}{
\begin{aligned}
\sqrt{v_1} &= \alpha+\alpha_1-\alpha_2-\alpha_3\\
\sqrt{v_2} &= \alpha-\alpha_1+\alpha_2-\alpha_3\\
\sqrt{v_3} &= \alpha-\alpha_1-\alpha_2+\alpha_3\\
a_1 &= \alpha+\alpha_1+\alpha_2+\alpha_3.
\end{aligned}
}
From this it then follows:
\weq{15}{
\begin{aligned}
4\alpha &= a_1+\sqrt{v_1}+\sqrt{v_2}+\sqrt{v_3}\\
4\alpha_1 &= a_1+\sqrt{v_1}-\sqrt{v_2}-\sqrt{v_3}\\
4\alpha_2 &= a_1-\sqrt{v_1}+\sqrt{v_2}-\sqrt{v_3}\\
4\alpha_3 &= a_1-\sqrt{v_1}-\sqrt{v_2}+\sqrt{v_3}.
\end{aligned}
}
If one gives $\sqrt{D}$ in (8) the opposite sign, and replaces $y$ by $-y$, then by (10) one must interchange $y_1$, $y_2$ with $-y_2$, $-y_1$. Thus $v_1$ remains unchanged, and $v_2$ and $v_3$ are interchanged with one another. If one takes for $y$ another root of the equation (8), say $y_1$, then on account of (9) $y_1$, $y_2$ must be replaced by $y_2$, $y$, and the quantities $v_1$, $v_2$, $v_3$ undergo a cyclic permutation, whereby the quantities $\alpha_1$, $\alpha_2$, $\alpha_3$ in (15) are also cyclically permuted. If, finally, one chooses the signs of $\sqrt{v_1}$, $\sqrt{v_2}$, $\sqrt{v_3}$ differently, but in such a way that (13) always remains satisfied, then certain permutations likewise occur among the quantities (15).
